\hypertarget{a-criticism-of-ghost}{%
\subsubsection{A Criticism of GHOST}\label{a-criticism-of-ghost}}

GHOST allows for blocks to link to a single canonical parent and multiple \emph{uncle} blocks.
In the full GHOST algorithm, uncle blocks contribute \emph{weight} to the canonical chain-segment, but do not contribute \emph{transactions}.
Thus, uncles have \emph{no ability} to substantially contribute to the canonical chain's \emph{state}.

Consider an empty-block DoS against a chain using GHOST.
If an attacker were to perform an empty-block DoS, the attacker could link back to honest miners' blocks as uncles, but never parents.
Given this, there is no easy way for the honest miners to end or mitigate the DoS.
The attacker can include honest miners' chain-work in a purely
\emph{beneficial} way --- there is a symmetry, thus honest miners (and the network) are at the mercy of the attacker.

Why does this symmetry exist? Because the \emph{cumulative weight} of each block (including uncles) is \emph{divorced} from \emph{the set of transactions} that is contributed by that block.
However, with a full block-DAG, when an attacker links to uncles in this way \emph{they must allow for the execution of all non-conflicting transactions} (i.e., those which would not cause a doublespend to occur).
Thus, GHOST
\emph{does not mitigate} empty-block DoS attacks; \emph{only} a full block-DAG can do that.
